% TOP_PROBLEM: {"id": 30006854, "title": "Calderón uniqueness for bounded measurable real isotropic conductivities", "category_id": 9, "category": {"id": 9, "name": "pde", "display_name": "Partial Differential Equations", "description": "PDEs and their applications in physics and geometry.", "slug": "pde", "order_index": 9, "created_at": "2026-07-31T15:26:25.671Z"}, "status": "open"}
% FIRST_POSTED: 2026-10-06
% ATTEMPT_STATUS: solved
% OPENAI_CLAIM: {"schema_version":1,"source_repo":"https://github.com/openai/math","source_commit":"adc7f1241b42e322a6451854ab7e4b4c146bf78a","release_date":"2026-10-06","match":"full","resolution":"disproved","independently_reviewed":false,"second_pass":"agreed","families":[{"family":"365","title":"Joint metric and connection recovery from one boundary patch","subject":"Partial differential equations","manuscripts":[{"title":"Nonuniqueness for Bounded Measurable Scalar Conductivities in Three Dimensions","pdf_path":"preprints/Nonuniqueness-for-Bounded-Measurable-Scalar-Conductivities-in-Three-Dimensions-September-23-2026/paper.pdf","readme_path":"preprints/Nonuniqueness-for-Bounded-Measurable-Scalar-Conductivities-in-Three-Dimensions-September-23-2026/README.md","date":"2026-09-23","theorem_ref":"thm:main (paper.tex:79)","statement_sources":[{"path":"preprints/Nonuniqueness-for-Bounded-Measurable-Scalar-Conductivities-in-Three-Dimensions-September-23-2026/build/paper.tex","line":79,"label":"thm:main"}]},{"title":"Smooth anisotropic uniqueness in the Calderón problem from one boundary patch","pdf_path":"preprints/Smooth-Anisotropic-Uniqueness-in-the-Calderon-Problem-from-One-Boundary-Patch-September-24-2026/paper.pdf","readme_path":"preprints/Smooth-Anisotropic-Uniqueness-in-the-Calderon-Problem-from-One-Boundary-Patch-September-24-2026/README.md","date":"2026-09-24","theorem_ref":"main:full (sections/01-introduction.tex:84); main:patch (sections/01-introduction.tex:54)","statement_sources":[{"path":"preprints/Smooth-Anisotropic-Uniqueness-in-the-Calderon-Problem-from-One-Boundary-Patch-September-24-2026/build/sections/01-introduction.tex","line":84,"label":"main:full"},{"path":"preprints/Smooth-Anisotropic-Uniqueness-in-the-Calderon-Problem-from-One-Boundary-Patch-September-24-2026/build/sections/01-introduction.tex","line":54,"label":"main:patch"}]}],"lean":{"doc_path":"lean/docs/365.md","comparators":["lean/ComparatorChallenges/Conductivity.json","lean/ComparatorChallenges/Conductivity.lean"],"declarations":["OAI.ScalarConductivity.main_nonuniqueness"],"covers_main_theorem":false},"reasoning_trace":null}]}
\documentclass[11pt]{article}
\usepackage[margin=25mm]{geometry}
\usepackage{fontspec}
\setmainfont{Latin Modern Roman}
\usepackage{amsmath,amssymb,mathtools}
\usepackage{unicode-math}
\setmathfont{Latin Modern Math}
\usepackage{enumitem}
\usepackage{xurl,hyperref,bookmark}
\hypersetup{colorlinks=true,urlcolor=blue}
\setcounter{secnumdepth}{0}
\setlength{\parindent}{0pt}
\setlength{\parskip}{5pt}
\setlength{\emergencystretch}{3em}
\begin{document}
\section*{Calderón uniqueness for bounded measurable real isotropic conductivities}

\subsection{OpenAI's claim}
\textit{Measurable scalar conductivity nonuniqueness}: The paper claims distinct bounded uniformly positive scalar conductivities in a three-dimensional ball with identical full Dirichlet-to-Neumann maps; both agree with the constant conductivity near the boundary.

\textit{One-patch anisotropic uniqueness}: The paper claims determination of a smooth Riemannian metric in dimension at least three from its zero-frequency energy form on any one open boundary patch; the only ambiguity is a diffeomorphism equal to the identity on that patch.

\subsection{Mathematical statement of the claimed result}
\paragraph{Mathematical claim 1.} In the three-dimensional ball $\Omega=B(0,3)$ there are distinct real scalar conductivities $\gamma_0,\gamma_1\in L^\infty(\Omega)$ with $0<c\le\gamma_j\le C<\infty$ almost everywhere, both equal to $1$ near the boundary, such that \[|\{\gamma_0\ne\gamma_1\}|>0,\qquad\Lambda_{\gamma_0}=\Lambda_{\gamma_1}:H^{1/2}(\partial\Omega)\to H^{-1/2}(\partial\Omega).\] Here $\Lambda_\gamma$ is the weak Dirichlet-to-Neumann operator for $\nabla\cdot(\gamma\nabla u)=0$. The equality is full-boundary equality despite positive-measure disagreement of bounded uniformly elliptic isotropic coefficients.

\paragraph{Mathematical claim 2.} Fix a dimension $n\ge3$. Assume the smooth $n$-manifold $M$ is compact and connected, with smooth nonempty boundary. For a proper relatively open observation patch $\Gamma\ne\varnothing$, the claim asserts \[\Lambda_{g_1,\Gamma}=\Lambda_{g_2,\Gamma}\quad\Longrightarrow\quad g_2=\Phi^*g_1,\qquad\Phi|_\Gamma=\operatorname{id},\] for some diffeomorphism smooth through the boundary. Both metrics are smooth, and the local maps are the harmonic Dirichlet energy forms for data supported in $\Gamma$. Equality for all full-boundary data gives a diffeomorphism fixing the entire boundary. This concerns smooth anisotropic metrics, not arbitrary bounded measurable scalar conductivities.

\subsection{Outline of the argument}
\textit{Measurable scalar conductivity nonuniqueness}: The construction first uses matrix-conductivity geometry to create a boundary-response obstruction. Scalarization preserves exact responses for a prescribed finite set of inputs, with the remaining inputs added through a separate compatibility argument. The limit retains uniform ellipticity and exact full-map equality rather than only a homogenized approximation.

\textit{One-patch anisotropic uniqueness}: Boundary jets and a common exterior cap identify actual Green kernels in an initial region. Remote source solutions produce harmonic coordinates and enough second-order jets. A mixed Green-kernel continuation theorem extends metric equality across a smooth frontier, so the maximal paired coordinate region cannot stop in the interior.

\subsection{Scope relative to this problem}
\textit{Measurable scalar conductivity nonuniqueness}: The constructed pair consists of distinct real bounded scalar coefficients, both measurable and uniformly elliptic, on a smooth three-dimensional ball, equal near its boundary, with identical full weak Dirichlet-to-Neumann operators on the required trace spaces. This admissible example refutes the universal n>=3 scalar uniqueness claim.

\textit{One-patch anisotropic uniqueness}: Corollary main:scalar proves isotropic uniqueness for smooth strictly positive scalar conductivities on smooth domains, even from one patch. The site requires arbitrary bounded measurable conductivities on Lipschitz domains, which includes nonsmooth classes absent from this theorem.

The claimed counterexample lies within the question’s hypotheses and contradicts its universal conclusion.

\subsection{Formal verification}
Family 365 describes Lean coverage for the measurable scalar-conductivity obstruction; smooth anisotropic and coupled connection recovery are separate.

\textit{Measurable scalar conductivity nonuniqueness}: The manuscript links \texttt{Conductivity.lean} (declaration \texttt{OAI.ScalarConductivity.main\_nonuniqueness}).

The documented selection includes this manuscript’s main assertion; that assertion retains the restrictions stated above.

\textit{One-patch anisotropic uniqueness}: No manuscript-specific Comparator is listed in the pinned inventory.

The documented selection does not cover the entire manuscript assertion relevant here.

This site has not run the Lean development. These links describe OpenAI’s selected formal statements and should be read with their stated scope.

\subsection{Status caveat}
OpenAI reports that an internal model produced the manuscripts. Their arguments have not been independently verified as proofs by this site. Any portion outside the documented formal scope may contain mistakes, and the catalogue label reports the announced mathematical claim.

\subsection{Sources}
\begin{itemize}
\item OpenAI, \emph{Nonuniqueness for Bounded Measurable Scalar Conductivities in Three Dimensions}, 2026-09-23; citation key \texttt{OAI:Nonuniqueness-for-Bounded-Measurable-Scalar-Conductivities-in-Three-Dimensions-September-23-2026}. \href{https://github.com/openai/math/blob/adc7f1241b42e322a6451854ab7e4b4c146bf78a/preprints/Nonuniqueness-for-Bounded-Measurable-Scalar-Conductivities-in-Three-Dimensions-September-23-2026/paper.pdf}{PDF (GitHub)}; \href{https://github.com/openai/math/raw/adc7f1241b42e322a6451854ab7e4b4c146bf78a/preprints/Nonuniqueness-for-Bounded-Measurable-Scalar-Conductivities-in-Three-Dimensions-September-23-2026/paper.pdf}{PDF (direct)}; \href{https://github.com/openai/math/blob/adc7f1241b42e322a6451854ab7e4b4c146bf78a/preprints/Nonuniqueness-for-Bounded-Measurable-Scalar-Conductivities-in-Three-Dimensions-September-23-2026/README.md}{Manuscript page and citation}; \href{https://github.com/openai/math/blob/main/preprints/Nonuniqueness-for-Bounded-Measurable-Scalar-Conductivities-in-Three-Dimensions-September-23-2026/paper.pdf}{Latest version}; Statement source: \href{https://github.com/openai/math/blob/adc7f1241b42e322a6451854ab7e4b4c146bf78a/preprints/Nonuniqueness-for-Bounded-Measurable-Scalar-Conductivities-in-Three-Dimensions-September-23-2026/build/paper.tex\#L79}{thm:main}.
\item OpenAI, \emph{Smooth anisotropic uniqueness in the Calderón problem from one boundary patch}, 2026-09-24; citation key \texttt{OAI:Smooth-Anisotropic-Uniqueness-in-the-Calderon-Problem-from-One-Boundary-Patch-September-24-2026}. \href{https://github.com/openai/math/blob/adc7f1241b42e322a6451854ab7e4b4c146bf78a/preprints/Smooth-Anisotropic-Uniqueness-in-the-Calderon-Problem-from-One-Boundary-Patch-September-24-2026/paper.pdf}{PDF (GitHub)}; \href{https://github.com/openai/math/raw/adc7f1241b42e322a6451854ab7e4b4c146bf78a/preprints/Smooth-Anisotropic-Uniqueness-in-the-Calderon-Problem-from-One-Boundary-Patch-September-24-2026/paper.pdf}{PDF (direct)}; \href{https://github.com/openai/math/blob/adc7f1241b42e322a6451854ab7e4b4c146bf78a/preprints/Smooth-Anisotropic-Uniqueness-in-the-Calderon-Problem-from-One-Boundary-Patch-September-24-2026/README.md}{Manuscript page and citation}; \href{https://github.com/openai/math/blob/main/preprints/Smooth-Anisotropic-Uniqueness-in-the-Calderon-Problem-from-One-Boundary-Patch-September-24-2026/paper.pdf}{Latest version}; Statement source: \href{https://github.com/openai/math/blob/adc7f1241b42e322a6451854ab7e4b4c146bf78a/preprints/Smooth-Anisotropic-Uniqueness-in-the-Calderon-Problem-from-One-Boundary-Patch-September-24-2026/build/sections/01-introduction.tex\#L84}{main:full}; Statement source: \href{https://github.com/openai/math/blob/adc7f1241b42e322a6451854ab7e4b4c146bf78a/preprints/Smooth-Anisotropic-Uniqueness-in-the-Calderon-Problem-from-One-Boundary-Patch-September-24-2026/build/sections/01-introduction.tex\#L54}{main:patch}.
\item Family 365: \href{https://github.com/openai/math/blob/adc7f1241b42e322a6451854ab7e4b4c146bf78a/lean/docs/365.md}{Lean scope}.
\item \href{https://github.com/openai/math/blob/adc7f1241b42e322a6451854ab7e4b4c146bf78a/lean/ComparatorChallenges/Conductivity.json}{Comparator Conductivity.json}.
\item \href{https://github.com/openai/math/blob/adc7f1241b42e322a6451854ab7e4b4c146bf78a/lean/ComparatorChallenges/Conductivity.lean}{Comparator Conductivity.lean}.
\item \href{https://github.com/openai/math/blob/adc7f1241b42e322a6451854ab7e4b4c146bf78a/CONTENTS.md}{OpenAI catalogue entry}.
\end{itemize}
\end{document}
